\section{Exact confinement under concatenation}
\label{sec:confinement}

Two confinement results govern the section.  Classifying recovery systems by
their outer functional labels gives the exact equation escape cost
(Theorem~\ref{thm:ungated-ranked-confinement}).  Deleting dispensable external
coefficients instead gives minimal-support transfer under a target-touching
distance hypothesis (Theorem~\ref{thm:minimal-support-confinement}).  The latter
transfers availability and support-based allocations even when equations can
escape.  Finally, the full label-indexed costs pass to the next concatenation
level by an associative min--sum law.

\paragraph{Represented inner code.}
Let $L/\F_q$ be a finite extension and let
$\iota:L\xrightarrow{\sim} I\leq\F_q^E$ be an $\F_q$-linear isomorphism,
used as the inner encoder.
Assume $0<\dim I<|E|$.
Whenever the outer alphabet $L$ occurs below, $I$ denotes the represented
inner code $(I,\iota)$, not only its image subspace.  In particular,
$O\circ I$, $\Phi_I$, $\lambda_{I,T}$, $\mu_{I,P,T}$, and
$\Gamma_{j,T}(O,I)$ depend on $\iota$.  The earlier pair
$K_P\subseteq D_P$ depends only on the image code and target/helper split.
Proposition~\ref{prop:functional-label-separation} shows that this distinction
cannot be suppressed under concatenation.

For an $L$-linear outer code $O\leq L^N$, write
\[
 O\circ I=\{(\iota(a_1),\ldots,\iota(a_N)):(a_1,\ldots,a_N)\in O\}.
\]
Each block has a label map
\[
 \underbrace{\F_q^E}_{\text{inner coefficients}}
 \xrightarrow{\ \Phi_I\ }
 \underbrace{L^*}_{\text{functional label}},
\qquad
\Phi_I(y)(a)=\langle y,\iota(a)\rangle.
\]
Here $L^*=\operatorname{Hom}_{\F_q}(L,\F_q)$.  The chosen message coordinates
identify each earlier generator column with the functional
$a\mapsto\iota(a)_e$; equivalently, it is the unique trace coefficient $b_e$
with $\iota(a)_e=\operatorname{Tr}_{L/\F_q}(b_ea)$.  We transport $T$ and
$G_P\alpha=\operatorname{id}_T$ through this same identification.
Two coefficient vectors receive the same label precisely when they differ by
an inner dual word, because $\ker\Phi_I=I^\perp$.  The prescribed-coset costs
below record the cheapest coefficient support in each fibre of this map.

\paragraph{Outer compatibility.}
A block vector
$y=(y_1,\ldots,y_N)$ belongs to $(O\circ I)^\perp$ precisely when
\begin{equation}
 (\Phi_I(y_1),\ldots,\Phi_I(y_N))\in
 \operatorname{FD}(O),
 \label{eq:block-functional-dual}
\end{equation}
where
\[
 \operatorname{FD}(O)=
 \{(\beta_j)\in(L^*)^N:\textstyle\sum_j\beta_j(a_j)=0
       \text{ for every }(a_j)\in O\}.
\]
Indeed, both conditions say that $y$ is orthogonal to every encoded outer
word.  Write $\beta_j(a)=\operatorname{Tr}_{L/\F_q}(b_ja)$.  The functional-dual
condition becomes
\[
 \operatorname{Tr}_{L/\F_q}\!\left(\sum_j b_ja_j\right)=0
 \qquad\text{for every }(a_j)\in O.
\]
Because $O$ is $L$-linear, the same holds after multiplying $(a_j)$ by every
$\lambda\in L$.  Nondegeneracy of the trace pairing then forces
$\sum_jb_ja_j=0$.  Hence $(b_j)\in O^\perp$, and the converse is immediate.
This identification preserves block support, so the support distance of
$\operatorname{FD}(O)$ is $d(O^\perp)$.

\paragraph{Confinement and target normalization.}
Fix a target block and a nonzero target-message subspace $T\leq U_P$.  A
system is \emph{confined} if every one of its equations is supported in that
target block.  Let
$\rho_T(I)$ be the minimum helper-union size of a normalized recovery system
for $T$ in the inner code, with value $\infty$ when no such system exists.
Thus $\rho_T(I)<\infty$ exactly when $T\leq W_P$.  In a concatenation, a
normalized system for $T$ means a linear family of dual equations whose
target-block coefficient map $\alpha:T\to\F_q^P$ satisfies
$G_P\alpha=\operatorname{id}_T$ under the fixed inner-message
identification.  This fixes a normalization condition, not a particular
section $\alpha$: the admissible sections vary in the minimization.
This agrees with the intrinsic target image of the
concatenated generator whenever the outer projection onto the chosen block is
surjective.  In particular, $d(O^\perp)>1$ guarantees surjectivity: an
$L$-linear coordinate image is either $0$ or $L$, and a zero image would put a
weight-one word in $O^\perp$.

For an $\F_q$-linear map $B:T\to L^*$, define
\begin{equation}
 \lambda_{I,T}(B)=
 \min_{\substack{Y:T\to\F_q^E\ \mathrm{linear}\\
                  \Phi_IY=B}}
 |\supp Y(T)|.
 \label{eq:ordinary-joint-coset-cost}
\end{equation}
For the target block, define
\begin{equation}
 \mu_{I,P,T}(B)=
 \min_{\substack{\alpha:T\to\F_q^P,\ \beta:T\to\F_q^J\ \mathrm{linear}\\
                  G_P\alpha=\operatorname{id}_T\\
                  \Phi_I(\alpha,\beta)=B}}
 |\supp\beta(T)|.
 \label{eq:target-joint-coset-cost}
\end{equation}
An empty minimum is $\infty$.  After choosing a basis of $T$,
\eqref{eq:ordinary-joint-coset-cost} is the minimum union support of
representatives of prescribed cosets of $I^\perp$.  This is the fixed-instance
quantity underlying generalized coset weights and generalized covering radii
\cite{ZhangEtAl2019,ElimelechFirerSchwartz2021}.  The target version retains
the normalization required by the recovery problem.  In particular,
\begin{equation}
 \mu_{I,P,T}(0)=\rho_T(I).
 \label{eq:zero-target-joint-coset-cost}
\end{equation}
For each equation parameter $t\in T$, the value $B(t)$ is exactly the outer
functional label induced by that equation in the block.

Let
\[
 \mathcal B_T(O)=
 \operatorname{Hom}_{\F_q}\bigl(T,\operatorname{FD}(O)\bigr).
\]
For $B\in\mathcal B_T(O)$, write $B=(B_h)_{h=1}^N$ with
$B_h:T\to L^*$.  Assume $N\geq2$, fix a target block $j$ whose coordinate projection
$O\to L$ is nonzero, hence surjective, and put
\begin{equation}
 \begin{split}
 \Gamma_{j,T}(O,I)=\min\Biggl\{&
 \mu_{I,P,T}(0)+d(I^\perp),\\
 &\min_{0\ne B\in\mathcal B_T(O)}
 \left(
 \mu_{I,P,T}(B_j)+
 \sum_{h\ne j}\lambda_{I,T}(B_h)
 \right)
 \Biggr\}.
 \end{split}
 \label{eq:weighted-subspace-cost}
\end{equation}

\begin{theorem}[Exact nonconfinement cost from labelled coset supports]
\label{thm:ungated-ranked-confinement}
\coverage{absent}%
\imports{Zhang:tuple-coset-support}%
Under the hypotheses above, $\Gamma_{j,T}(O,I)$ is the least helper-union
size of a nonconfined normalized recovery system for $T$ in $O\circ I$, with
value $\infty$ if no such system exists.
Thus every cost-$\leq r$ system for $T$ is confined if and only if
\begin{equation}
 r<\Gamma_{j,T}(O,I).
 \label{eq:weighted-subspace-gate}
\end{equation}
For $1\leq t\leq\ell=\dim W_P$, set
\begin{equation}
 \Gamma_{j,t}(O,I)=
 \min_{\substack{T\leq W_P\\\dim T=t}}\Gamma_{j,T}(O,I).
 \label{eq:ranked-weighted-cost}
\end{equation}
Then every cost-$\leq r$ system for every helper-recoverable
$t$-dimensional target subspace is confined if and only if
$r<\Gamma_{j,t}(O,I)$.  Below either threshold, restriction and zero-extension
are inverse bijections on normalized systems and preserve exact helper
supports.
\end{theorem}

\begin{proof}
Write a recovery system as an $\F_q$-linear map
$Y:T\to(O\circ I)^\perp$ and let $Y_h$ be its block maps.  By
\eqref{eq:block-functional-dual},
\[
 B_h=\Phi_IY_h
\]
defines a linear map $B:T\to\operatorname{FD}(O)$.  Conversely, compatible
linear lifts of the block maps of any such $B$ give a map into the
concatenated dual.  At the target block, compatibility with the chosen target
space is exactly the constraint in~\eqref{eq:target-joint-coset-cost}.

Suppose first that $B\ne0$.  No nonzero element of
$\operatorname{FD}(O)$ is supported only at $j$: testing such an element
against outer words whose $j$th coordinate ranges over $L$ would force its
$j$th functional to vanish.  Hence some block map $B_h$ with $h\ne j$ is
nonzero, and every lift is nonconfined.  Distinct inner blocks have disjoint
helper coordinates, so the minimum helper union for this fixed $B$ is
\[
 \mu_{I,P,T}(B_j)+
 \sum_{h\ne j}\lambda_{I,T}(B_h).
\]
The block lifts can be chosen independently, so this lower bound is attained.

If $B=0$, every block map takes values in $I^\perp=\ker\Phi_I$.  The target
block costs $\rho_T(I)=\mu_{I,P,T}(0)$.  Nonconfinement requires a nonzero map
$T\to I^\perp$ in another block.  Its image has union support at least
$d(I^\perp)$, and equality is attained by $u\mapsto f(u)v$, where
$0\ne f\in T^*$ and $v$ is a minimum-weight word of $I^\perp$.

The zero and nonzero functional sectors are exhaustive, proving
\eqref{eq:weighted-subspace-cost} and~\eqref{eq:weighted-subspace-gate}.
Minimizing over $t$-dimensional $T$ proves the ranked statement.  Below the
first nonconfined cost, every global system has zero coefficient maps in all
blocks outside $j$.  Restriction therefore preserves its target normalization
and helper support.  Conversely, zero-extending an inner dual system remains
orthogonal to every concatenated word, so the two operations are inverse with
coefficients and supports unchanged.
\proves{thm:ungated-ranked-confinement}%
\end{proof}

\paragraph{The running example: equations versus useful repairs.}
For the code in the introduction,
$G=(e_1\mid e_1,e_1,e_2,e_1+e_2)$.
Its inner-dual distance is two: two repeated columns give a weight-two
relation, and no column is zero.  The local recovery cost for $u$ is one.
Identify its message space with $L=\F_4$ and take $O=L^N$, $N\geq2$.
There are no nonzero outer-dual labels, so the exact formula gives
$\Gamma_{j,1}=1+2=3$.  A cost-three nonconfined equation adds the relation
$u+u=0$ in another block to the local recovery of $u$.
Deleting those two external coefficients leaves the original repair.

Nevertheless every minimal repair remains in its own block at every radius:
the outer messages are independent, so external coefficients cannot supply
any part of the target functional.  The three local choices remain
$\{1\}$, $\{2\}$, and $\{3,4\}$.  This example exhibits the zero sector
of the exact formula and the difference between its coefficient threshold
and the operational theorem below.  The nonzero-sector alignment obstruction
is exhibited separately in Proposition~\ref{prop:functional-label-separation}.

\subsection{Minimal supports and operational transfer}

An external equation can violate coefficient confinement without supplying a
new recovery option.  To separate these questions, let
$\mathcal M_{j,T}^{(\leq r)}(O\circ I)$ be the inclusion-minimal helper
supports recovering $T$ with at most $r$ helpers.  A set recovers $T$ when
some normalized system uses only that set.  Define the target-touching distance
\[
 \delta_j(O^\perp)=\min\{\wt(u):u\in O^\perp,\ u_j\ne0\},
\]
with empty minimum $\infty$.

\begin{theorem}[Minimal-support confinement]
\label{thm:minimal-support-confinement}
\coverage{absent}%
Let $O\leq L^N$ be $L$-linear, let its projection onto block $j$ be
surjective, and let $0\ne T\leq U_P$.  If $\delta_j(O^\perp)>r+1$, then
\[
 \mathcal M_{j,T}^{(\leq r)}(O\circ I)
 =\{\{j\}\times H:H\in\mathcal M_T^{(\leq r)}(I)\}.
\]
No inequality involving $\rho_T(I)+d(I^\perp)$ is required.
Consequently bounded recovery success has the same probability for every
global helper-availability law with the prescribed target-block marginal.
\end{theorem}

\begin{proof}
Let $Y$ be a normalized system with helper support $H$, $|H|\leq r$, and
put $B_h=\Phi_IY_h$.  If $B_j\ne0$, choose $v\in T$ with $B_j(v)\ne0$.
The trace identification sends $B(v)$ to an outer dual word touching $j$.
Its other nonzero blocks each contain a helper in $H$, so its weight is at
most $r+1$, a contradiction.  Thus $B_j=0$, and the target-block portion
of $Y$ is an inner-dual system with the same target normalization.  Deleting
all external coefficients preserves recovery and only decreases support.
Every global minimal support is therefore local.  Conversely, a proper
globally recovering subset of a local minimal support would, by this same
restriction argument, contradict local minimality.  Zero-extension supplies
the other inclusion.  For each available set, success means containing one
of these minimal supports; the two success indicators therefore agree
pointwise after restriction to the target block.  Taking expectations proves
the probability assertion without an independence assumption.
\proves{thm:minimal-support-confinement}%
\end{proof}

The first new external repair has its own threshold.  Let
$\varepsilon_{j,T}$ be the least size of a globally recovering helper set
$H$ whose portion in block $j$ does not recover $T$ locally.  Equivalently,
it is the least size of a global minimal support using an external helper:
minimize inside $H$ for one direction; a local recovery inside a minimal
support contradicts minimality for the other.  The theorem gives
$\varepsilon_{j,T}\geq\delta_j(O^\perp)-1$ when the latter is finite,
and $\varepsilon_{j,T}=\infty$ when $\delta_j(O^\perp)=\infty$.
Its value cannot be obtained simply by deleting the zero sector from
$\Gamma_{j,T}$: a nonzero-labelled lift can also contain a local recovery.
For $O=L^N$, every minimal recovery support is local at every radius,
although adjoining an external inner-dual word gives the finite equation
threshold $\rho_T(I)+d(I^\perp)$ whenever $T$ is locally recoverable.

\subsection{Repeated concatenation}

The ordinary prescribed-coset costs retain the intermediate labels needed at
the next concatenation level.  Target-normalized composition also retains the
part of each intermediate label already supplied by target coordinates.  In
schematic form,
\[
 \text{cost of output label}
 =\min_{\text{compatible intermediate labels}}
   \sum_{\text{inner blocks}}\text{inner realization cost}.
\]
The eliminated variables are the intermediate labels shared by adjacent
levels.  The resulting formulas are instances of variable elimination over
the min-sum semiring \cite{AjiMcEliece2000}.  This is also the algebraic
form of passing symbol-indexed inner costs to an outer decoder, as in soft and
generalized-concatenated decoding \cite{GuruswamiSudan2002,BlomqvistEtAl2020}.

The zero-functional nonconfinement cost alone is not compositional.  The next
separation holds the underlying binary code and all its unlabelled recovery
data fixed.

\begin{proposition}[Unlabelled data do not determine finite composition]
\label{prop:functional-label-separation}
\coverage{absent}%
Over $\F_2\subseteq\F_4$, there are two represented inner encoders with the
same image code, dual code, target/helper nested pair, complete relative-weight
hierarchy, dual distance, unlabelled normalized helper-support family, and
zero-functional nonconfinement costs for every recoverable target subspace, but one
fixed outer represented code gives different exact nonconfinement costs.
\end{proposition}

\begin{proof}
Write the nonzero elements of $\F_4$ as
$1,\omega,\omega^2=1+\omega$ and take generator columns
\[
 I_1:(1,1,\omega),
 \qquad
 I_2:(1,1,\omega^2),
\]
with the first coordinate targeted.  Relative to the binary basis
$(1,\omega)$ of $\F_4$, each displayed field element is a two-dimensional
binary generator column.  Both generator matrices have binary row
space
\[
 \{(x,x,y):x,y\in\F_2\}
\]
and dual $\langle(1,1,0)\rangle$.  Thus both have dual distance $2$,
$K_P=0\subseteq D_P=\langle(1,0)\rangle$, $M_1=1$, and the same unique
minimal helper support.  Their recoverable target space is one-dimensional,
so their complete zero-functional profile is the common value $1+2=3$.

For this table identify a column $c$ with the functional
$a\mapsto\Tr(\omega^2ca)$.  This is the coordinate-dot-product pairing
in the basis $(1,\omega)$; the corresponding trace coefficient is
$\omega^2c$.  The common rescaling preserves every outer $L$-linear dual
subspace.  In this column-label order $(1,\omega,\omega^2)$, the ordinary coset costs are
$(1,1,2)$ for $I_1$ and $(1,2,1)$ for $I_2$; the normalized target costs are
respectively $(0,2,1)$ and $(0,1,2)$.  Take the fixed outer code whose
functional dual is the line spanned by $(1,\omega)$.  Its nonzero multiples
have target label $s$ and external label $s\omega$.  The complete calculation
is
\[
\begin{array}{c|cc|ccc|ccc}
 s & s & s\omega
   & \mu_1(s) & \lambda_1(s\omega) & \text{total}
   & \mu_2(s) & \lambda_2(s\omega) & \text{total}\\ \hline
 1        & 1        & \omega   & 0&1&1 & 0&2&2\\
 \omega   & \omega   & \omega^2 & 2&2&4 & 1&1&2\\
 \omega^2 & \omega^2 & 1        & 1&1&2 & 2&1&3
\end{array}
\]
Equation~\eqref{eq:weighted-subspace-cost} therefore gives exact costs $1$
and $2$.  Only the alignment of costs with the intermediate functional labels
has changed.
\uses{thm:ungated-ranked-confinement}%
\proves{prop:functional-label-separation}%
\end{proof}

Let $\F_q\subseteq L\subseteq M$ be a tower of finite fields.  Let $B$ be an
$\F_q$-linear represented code with message space $L$, and let $A$ be an
$L$-linear represented code with message space $M$.  The composite $A\circ B$
replaces each coordinate of an $A$-word by its represented $B$-block.
After identifying a functional with its unique trace coefficient, write their
dual restriction maps as
\[
 \phi_B:\F_q^{E_B}\longrightarrow L,
 \qquad
 \phi_A:L^{E_A}\longrightarrow M.
\]
For an $\F_q$-space $T$ and $b:T\to L$, put
\[
 \Lambda_{B,T}(b)=
 \min_{\phi_BY=b}|\supp Y(T)|.
\]
Define $\Lambda_{A,T}$ and $\Lambda_{A\circ B,T}$ similarly.  If
$J\subseteq E_B$, let $\Lambda_{B,J,T}$ denote the same minimum with $Y$
supported on $J$ and $\phi_B$ restricted to those coordinates.
Thus $\phi$ is the trace-coordinate form of $\Phi$, and $\Lambda$ is the
corresponding ordinary prescribed-coset cost $\lambda$.  The change of letters
only distinguishes the two levels in the following substitution.

For the recovery-specific form, suppose that a target set in $A\circ B$ meets
block $e$ in $P_e\subseteq E_B$.  Put $J_e=E_B\setminus P_e$ and
$U_{P_e}=\operatorname{im}\phi_{B,P_e}$.  If $T\leq M$, let
$\iota_T:T\hookrightarrow M$ be the inclusion.  The variables below have the
following roles:
\[
\begin{array}{c|l}
 U_e & \text{intermediate contribution supplied by target coordinates in block }e\\
 X_e & \text{total intermediate functional label required in block }e\\
 X_e-U_e & \text{label that the helper coordinates in block }e\text{ must realize.}
\end{array}
\]

\begin{theorem}[Min--sum closure of prescribed-coset support costs]
\label{prop:prescribed-coset-composition}
\coverage{absent}%
\imports{Zhang:tuple-coset-support,Aji:distributive-law}%
\uses{thm:ungated-ranked-confinement}%
For every $\F_q$-linear map $c:T\to M$, the ordinary prescribed-coset cost is
\begin{equation}
 \Lambda_{A\circ B,T}(c)=
 \min_{\substack{X:T\to L^{E_A}\ \phi_AX=c}}
 \sum_{e\in E_A}\Lambda_{B,T}(X_e).
 \label{eq:coset-cost-composition}
\end{equation}
For $T\leq M$ and the target split above, the target-normalized recovery cost
is
\begin{equation}
 \mu_{A\circ B,P,T}(c)=
 \min_{\substack{U,X:T\to L^{E_A}\\
                  U_e(T)\subseteq U_{P_e}\ (e\in E_A)\\
                  \phi_AU=\iota_T,\ \phi_AX=c}}
 \sum_{e\in E_A}\Lambda_{B,J_e,T}(X_e-U_e).
 \label{eq:target-cost-composition}
\end{equation}
Both formulas iterate associatively through a finite tower: either
parenthesization minimizes over the same intermediate labels and target
contributions and sums the same leaf costs.  Storing minimizing target and
helper lifts propagates one coefficient-level witness; retaining every witness
requires the full lift sets rather than only the numerical minima.
\end{theorem}

\begin{proof}
A leaf-level lift $Y:T\to\F_q^{E_A\times E_B}$ determines intermediate maps
$X_e=\phi_BY_e$.  Its support is the disjoint union of its supports in the
inner blocks.  At fixed $X$, the block lifts are independent and attain their
minima separately, which proves \eqref{eq:coset-cost-composition}.

For the target-normalized formula, lift $U_e$ to an actual target map
$a_e:T\to\F_q^{P_e}$.  The condition $\phi_AU=\iota_T$ is exactly target
normalization.  Once $X_e$ is fixed, the helper contribution must lift
$X_e-U_e$ through $\phi_{B,J_e}$, independently in each block.  This proves
\eqref{eq:target-cost-composition}; storing the displayed lifts gives the
stated witness semantics.  At three or more levels, either parenthesization
eliminates the same intermediate arrays and target contributions, proving
associativity.
\uses{thm:ungated-ranked-confinement}%
\proves{prop:prescribed-coset-composition}%
\end{proof}

\paragraph{Sharp scalar envelopes.}
The full labelled function gives simple bounds when only its smallest and
largest nonzero values are retained.  When $T\ne0$, set
\[
 \delta_{B,T}=\min_{0\ne b:T\to L}\Lambda_{B,T}(b),
 \qquad
 R_{B,T}=\max_{0\ne b:T\to L}\Lambda_{B,T}(b).
\]
For $c\ne0$, the exact formula has the sharp coarsening
\begin{equation}
 \delta_{B,T}\Lambda_{A,T}(c)
 \leq \Lambda_{A\circ B,T}(c)
 \leq R_{B,T}\Lambda_{A,T}(c).
 \label{eq:coset-cost-envelope}
\end{equation}
Neither constant can be improved without further hypotheses.
Indeed, every nonzero intermediate coordinate map costs between
$\delta_{B,T}$ and $R_{B,T}$.  Minimizing the number of nonzero intermediate
coordinates proves the bounds; a one-coordinate outer realization attaining
either extremum proves sharpness.

\paragraph{Dual-distance recursion.}
The same zero/nonzero-label split gives
\begin{equation}
 d((A\circ B)^\perp)=
 \min\left\{
 d(B^\perp),
 \min_{0\ne z\in A^\perp}
 \sum_{e\in E_A}\Lambda_{B,\F_q}(u\mapsto uz_e)
 \right\}.
 \label{eq:dual-distance-composition}
\end{equation}
A word in $(A\circ B)^\perp$ either has zero intermediate label, when a
minimum word of $B^\perp$ in one block is optimal, or has a nonzero label
$z\in A^\perp$, when the block representatives minimize independently.  This
proves \eqref{eq:dual-distance-composition}.

\paragraph{Tower consequence.}
Substituting these recursively computed costs into
\eqref{eq:weighted-subspace-cost} gives the exact nonconfinement cost through
any finite tower, with the zero and nonzero functional sectors kept separate
at every level.

The trace identifications in these formulas are compatible with the tower.
Write $\iota_A:M\to A\leq L^{E_A}$ for the represented encoder of $A$.
For $v\in T$ and $m\in M$, trace transitivity gives
\[
 \sum_e\Tr_{L/\F_q}\bigl(X_e(v)\,\iota_A(m)_e\bigr)
 =\Tr_{M/\F_q}\bigl(\phi_A(X(v))m\bigr),
\]
and nondegeneracy of the trace pairing identifies the composite restriction
map with $\phi_A\circ\phi_B^{E_A}$, where $\phi_B^{E_A}$ applies $\phi_B$
blockwise.

\paragraph{Simultaneous confinement in every dimension.}
\begin{corollary}[Rank-one criterion for complete bounded transfer]
\label{cor:all-rank-bottleneck}
\coverage{absent}%
For every target block $j$ with nonzero outer projection,
\begin{equation}
 \min_{0\ne T\leq W_P}\Gamma_{j,T}(O,I)=\Gamma_{j,1}(O,I).
 \label{eq:all-rank-bottleneck}
\end{equation}
For every integer $r\geq0$, the following are equivalent:
\begin{enumerate}
 \item every cost-$\leq r$ system on every helper-recoverable target line
 is confined;
 \item every cost-$\leq r$ system on every nonzero helper-recoverable
 target subspace is confined; and
 \item $r<\Gamma_{j,1}(O,I)$.
\end{enumerate}
Under these conditions, restriction and zero-extension are inverse bijections
simultaneously for every nonzero $T\leq W_P$, preserving every coefficient and
exact helper support through radius $r$.

Consequently, any statistic defined solely from these bounded systems and
invariant under this bijection is preserved.  In particular, the bounded
reliability laws and the capacity and scheduling regions defined later from
their coefficients or supports agree.  Targetwise and rankwise minimum costs
also agree after truncation $c\mapsto\min\{c,r+1\}$ (with
$\min\{\infty,r+1\}=r+1$).  If every relevant inner minimum is at most $r$,
then the untruncated minima agree as well.
\end{corollary}

\begin{proof}
If a system on $0\ne T\leq W_P$ is nonconfined, some external block map is
nonzero on a vector $u\in T$.  Restriction to $\langle u\rangle$ remains
nonconfined and cannot enlarge its helper union.  This proves the displayed
equality, and Theorem~\ref{thm:ungated-ranked-confinement} gives the three
equivalent conditions.  When they hold, every bounded system is confined, so
restriction deletes only zero external blocks and zero-extension is its
inverse.  The coefficient, support, reliability, capacity, and scheduling
assertions follow from this bijection.  For a minimum helper-union cost, either
the inner minimum is at most $r$ and the bijection gives equality, or both
inner and concatenated minima exceed $r$; this is precisely equality after
truncation at $r+1$.
\uses{thm:ungated-ranked-confinement}%
\proves{cor:all-rank-bottleneck}%
\end{proof}
